\documentclass[../../main.tex]{subfiles}% 注意这里的文件路径不能用 ./main.tex ,否则用latexmk编译子文件会报错
\graphicspath{{\subfix{./image/}}{\subfix{../../image/}}} % 指定图片引用路径,后续可以直接使用图片文件名
% 如果图片存放在上级目录,则使用 ../../image/ 作为备用路径

% 例如:
% \begin{figure}[H]
% \centering
% \includegraphics[width=0.3\textwidth]{图.png}
% \caption{}
% \label{figure:图}
% \end{figure}
% 注意:上述\label{}一定要放在\caption{}之后,否则引用图片序号会只会显示??.

\begin{document}

\section{内积空间}

\subsection{定义与基本性质}

\begin{definition}
线性空间 $\mathscr{X}$ 上的一个二元函数 $a(\cdot,\cdot):\mathscr{X}\times\mathscr{X}\to\mathbb{K}$，称为是\textbf{共轭双线性函数}，如果
\begin{enumerate}[(1)]
\item $a(x,\alpha_1y_1+\alpha_2y_2)=\overline{\alpha_1}a(x,y_1)+\overline{\alpha_2}a(x,y_2)$,
\item $a(\alpha_1x_1+\alpha_2x_2,y)=\alpha_1a(x_1,y)+\alpha_2a(x_2,y)$,
\end{enumerate}
其中 $\forall x,y,x_1,x_2,y_1,y_2\in\mathscr{X},\forall\alpha_1,\alpha_2\in\mathbb{K}$. 称由
\begin{align*}
q(x)\triangleq a(x,x)\quad(\forall x\in\mathscr{X})
\end{align*}
定义的函数为 $\mathscr{X}$ 上\textbf{由 $\bm{a}$ 诱导的二次型}.
\end{definition}

\begin{proposition}\label{proposition:共轭双线性函数基本性质}
设 $a$ 是 $\mathscr{X}$ 上的共轭双线性函数，$q$ 是由 $a$ 诱导的二次型，那么
\begin{align*}
q(x)\in\mathbb{R}\ (\forall x\in\mathscr{X})\iff a(x,y)=\overline{a(y,x)}\quad(\forall x,y\in\mathscr{X}).
\end{align*}
\end{proposition}
\begin{proof}
充分性是显然的,下证必要性. 从
\begin{align*}
q(x+y)=\overline{q(x+y)}\quad(\forall x,y\in\mathscr{X}),
\end{align*}
可推出
\begin{align}
a(x,y)+a(y,x)=\overline{a(x,y)}+\overline{a(y,x)}.\label{eq:::--fiojoi3048HUAAIJOW1}
\end{align}
在 \eqref{eq:::--fiojoi3048HUAAIJOW1} 式中换 $y$ 为 $\mathrm{i}y$，即得
\begin{align}
-a(x,y)+a(y,x)=\overline{a(x,y)}-\overline{a(y,x)}.\label{eq:::--fiojoi3048HUAAIJOW2}
\end{align}
\eqref{eq:::--fiojoi3048HUAAIJOW1} 式与 \eqref{eq:::--fiojoi3048HUAAIJOW2} 式相减即得 $a(x,y)=\overline{a(y,x)}$.
\end{proof}

\begin{definition}
\begin{enumerate}
\item 线性空间 $\mathscr{X}$ 上的一个共轭双线性函数 $(\cdot,\cdot):\mathscr{X}\times\mathscr{X}\to\mathbb{K}$ 称为是一个\textbf{半内积}，如果它满足:
\begin{enumerate}[(1)]
\item 共轭对称性:$(x,y)=\overline{(y,x)}\quad(\forall x,y\in\mathscr{X})$;
\item 非负定性:$(x,x)\geqslant 0\ (\forall x\in\mathscr{X})$ .
\end{enumerate}
具有内积的线性空间称为\textbf{半内积空间}.

\item 线性空间 $\mathscr{X}$ 上的一个共轭双线性函数 $(\cdot,\cdot):\mathscr{X}\times\mathscr{X}\to\mathbb{K}$ 称为是一个\textbf{内积}，如果它满足:
\begin{enumerate}[(1)]
\item 共轭对称性:$(x,y)=\overline{(y,x)}\quad(\forall x,y\in\mathscr{X})$;
\item 正定性:$(x,x)\geqslant 0\quad(\forall x\in\mathscr{X})$, $(x,x)=0\iff x=\theta$ .
\end{enumerate}
具有内积的线性空间称为\textbf{内积空间}，记作 $(\mathscr{X},(\cdot,\cdot))$.
\end{enumerate}
\end{definition}

\begin{theorem}[Cauchy-Schwarz不等式]\label{theorem:泛函分析--Cauchy-Schwarz不等式}
设 $a$ 是线性空间 $\mathscr{X}$ 上的共轭双线性函数，$q(x)$ 是由 $a$ 诱导的二次型. 如果
\begin{align*}
q(x)\geqslant 0\quad(\forall x\in\mathscr{X})\quad\text{且}\quad q(x)=0\iff x=\theta,
\end{align*}
那么
\begin{align}
|a(x,y)|\leqslant [q(x)q(y)]^\frac{1}{2}\quad(\forall x,y\in\mathscr{X}),\label{eq:::--fiojoi3048HUAAIJOW5}
\end{align}
而且上式中的等号成立当且仅当 $x$ 与 $y$ 线性相关时成立.

特别地,若 $(\mathscr{X},(\cdot,\cdot))$ 是内积空间,令
\begin{align*}
\|x\|=(x,x)^\frac{1}{2}\quad(\forall x\in\mathscr{X}),
\end{align*}
则有
\begin{align*}
|(x,y)|\leqslant\|x\|\cdot\|y\|\quad(\forall x,y\in\mathscr{X}),
\end{align*}
而且上式中的等号成立当且仅当 $x$ 与 $y$ 线性相关时成立.
\end{theorem}
\begin{proof}
不妨设 $y\neq\theta$, 对 $\forall\lambda\in\mathbb{K}$ 考察
\begin{align}
q(x+\lambda y)=q(x)+\overline{\lambda}a(x,y)+\lambda a(y,x)+|\lambda|^2q(y)\geqslant 0.\label{eq:::--fiojoi3048HUAAIJOW6}
\end{align}
取 $\lambda=-\frac{a(x,y)}{q(y)}$, 由\refpro{proposition:共轭双线性函数基本性质}知$a(x,y)=\overline{a(y,x)}$, 所以
\begin{align*}
q(x)-\frac{2|a(x,y)|^2}{q(y)}+\frac{|a(x,y)|^2}{q(y)}\geqslant 0,
\end{align*}
由此立得 \eqref{eq:::--fiojoi3048HUAAIJOW5} 式. 又当 $x=-\lambda y\ (\lambda\in\mathbb{K})$ 时, \eqref{eq:::--fiojoi3048HUAAIJOW5} 式中的等号成立. 反之, 若 \eqref{eq:::--fiojoi3048HUAAIJOW5} 式中的等号成立, 则 \eqref{eq:::--fiojoi3048HUAAIJOW6} 式中的等号成立, 从而由条件可知$x=-\lambda y\ (\lambda\in\mathbb{K})$.

特别地,若 $(\mathscr{X},(\cdot,\cdot))$ 是内积空间,令
\begin{align*}
\|x\|=(x,x)^\frac{1}{2}\quad(\forall x\in\mathscr{X}),
\end{align*}
则取$\mathscr{X}$上的共轭双线性函数$a(x,y)=(x,y)$,由内积定义知,此时由$a$诱导的二次性$q(x)$也满足
\begin{align*}
q(x)\geqslant 0\quad(\forall x\in\mathscr{X})\quad\text{且}\quad q(x)=0\iff x=\theta.
\end{align*}
故由\eqref{eq:::--fiojoi3048HUAAIJOW5}式知
\begin{align*}
|(x,y)|\leqslant\|x\|\cdot\|y\|\quad(\forall x,y\in\mathscr{X}),
\end{align*}
而且上式中的等号成立当且仅当 $x$ 与 $y$ 线性相关时成立.
\end{proof}

\begin{proposition}[内积的连续性]\label{proposition:内积的连续性}
内积空间 $(\mathscr{X},(\cdot,\cdot))$ 按
\begin{align}\label{eq:::--fiojoi3048HUAAIJOW3}
\|x\|=(x,x)^\frac{1}{2}\quad(\forall x\in\mathscr{X})
\end{align}
定义范数$\|\cdot\|$,则$(\mathscr{X},\|\cdot\|)$是严格凸的 $B^*$ 空间.并且内积 $(x,y)$ 是 $\mathscr{X}\times\mathscr{X}$ 上关于范数 $\|\cdot\|$ 的连续函数.
而且还有
\begin{align*}
\left\| x+y \right\| ^2=\left\| x \right\| ^2+\left\| y \right\| ^2+2\mathrm{Re}\left( x,y \right) .
\end{align*}
\end{proposition}
\begin{proof}
只要证明按 \eqref{eq:::--fiojoi3048HUAAIJOW3} 式定义的 $\|\cdot\|$ 是范数. 事实上, 正定性和齐次性都是显然成立的, 下面来验证三角不等式.
\begin{align*}
\|x+y\|^2&=(x+y,x+y)
=(x,x)+(x,y)+(y,x)+(y,y)\\
&\leqslant\|x\|^2+2\|x\|\|y\|+\|y\|^2
=(\|x\|+\|y\|)^2\quad(\forall x,y\in\mathscr{X}).
\end{align*}

对$\forall 0<\lambda<1$, 根据\hyperref[theorem:泛函分析--Cauchy-Schwarz不等式]{Cauchy-Schwarz不等式},当$\|x\|=\|y\|=1,x\neq y$时,有
\begin{align*}
\parallel \lambda x+(1-\lambda )y\parallel ^2=\left( \lambda x+(1-\lambda )y,\lambda x+(1-\lambda )y \right) =\lambda ^2\parallel x\parallel ^2+2\lambda (1-\lambda )\mathrm{Re}(x,y)+(1-\lambda )^2\parallel y\parallel ^2<[\lambda +(1-\lambda )]^2=1.
\end{align*}

再证明内积的连续性.设 $x_n\to x,y_n\to y$. 那么由\refthe{theorem:度量空间中收敛序列极限的唯一性和序列有界性}知$\|x_n\|$ 和 $\|y_n\|$ 有界, 用 $M$ 表示它们的一个上界, 便有
\begin{align*}
|(x_n,y_n)-(x,y)|
&\leqslant|(x_n,y_n)-(x,y_n)|+|(x,y_n)-(x,y)|\\
&\leqslant\|x_n-x\|\cdot\|y_n\|+\|x\|\cdot\|y_n-y\|\\
&\leqslant M\|x_n-x\|+\|x\|\|y_n-y\|\to 0\quad(n\to\infty).
\end{align*}

再由内积的定义可得
\begin{align*}
\left\| x+y \right\| ^2&=\left( x+y,x+y \right) =\left\| x \right\| ^2+\left\| y \right\| ^2+\left( x,y \right) +\left( y,x \right) 
\\
&=\left\| x \right\| ^2+\left\| y \right\| ^2+\left( x,y \right) +\overline{\left( x,y \right)}
\\
&=\left\| x \right\| ^2+\left\| y \right\| ^2+2\mathrm{Re}\left( x,y \right) .
\end{align*}
\end{proof}

\begin{proposition}[平行四边形等式(极化恒等式)]\label{proposition:平行四边形等式--泛函分析}
在 $B^*$ 空间 $(\mathscr{X},\|\cdot\|)$ 中, 在 $\mathscr{X}$ 上可引入一个内积 $(\cdot,\cdot)$ 适合
\begin{align}
(x,x)^\frac{1}{2}=\|x\|\quad(\forall x\in\mathscr{X}),\label{eq:::--fiojoi3048HUAAIJOW7}
\end{align}
当且仅当范数 $\|\cdot\|$ 满足如下\textbf{平行四边形等式}(或\textbf{极化恒等式}):
\begin{align*}
\|x+y\|^2+\|x-y\|^2=2(\|x\|^2+\|y\|^2)\quad(\forall x,y\in\mathscr{X}).
\end{align*}
\end{proposition}
\begin{proof}
必要性可通过直接计算得到. 为了证充分性, 令
\begin{align*}
(x,y)=
\begin{cases}
\frac{1}{4}\left(\|x+y\|^2-\|x-y\|^2\right), & \text{当}\ \mathbb{K}=\mathbb{R}, \\
\frac{1}{4}\left(\|x+y\|^2-\|x-y\|^2+\mathrm{i}\|x+\mathrm{i}y\|^2-\mathrm{i}\|x-\mathrm{i}y\|^2\right), & \text{当}\ \mathbb{K}=\mathbb{C}.
\end{cases}
\end{align*}
容易验证它是一个满足 \eqref{eq:::--fiojoi3048HUAAIJOW7} 式的内积.
\end{proof}

\begin{definition}
完备的内积空间称为 $\mathbf{Hilbert}$\textbf{空间}.
\end{definition}

\begin{proposition}
\begin{enumerate}[(1)]
\item $\mathbb{R}^n,\mathbb{C}^n$ 都是Hilbert空间，它们的内积分别定义为
\begin{align*}
(x,y)=\sum\limits_{i=1}^n x_iy_i\quad(\forall x,y\in\mathbb{R}^n),
\quad
(x,y)=\sum\limits_{i=1}^n x_i\overline{y_i}\quad(\forall x,y\in\mathbb{C}^n),
\end{align*}
其中 $x=(x_1,x_2,\cdots,x_n),y=(y_1,y_2,\cdots,y_n)$.

\item $l^2$ 空间是Hilbert空间(定义见\refpro{proposition:常见B空间11111})，规定内积
\begin{align*}
(x,y)=\sum\limits_{i=1}^\infty x_i\overline{y_i},
\end{align*}
其中 $x=(x_1,x_2,\cdots,x_i,\cdots),y=(y_1,y_2,\cdots,y_i,\cdots)\in l^2$.

\item $L^2(\Omega,\mu)$ 是Hilbert空间(定义见\refpro{proposition:常见B空间11111})，规定内积
\begin{align*}
(u,v)=\int_\Omega u(x)\cdot\overline{v(x)}\mathrm{d}\mu,
\end{align*}
其中 $u,v\in L^2(\Omega,\mu)$.

\item 在空间 $C^k(\overline{\Omega})$ 中(定义见\refpro{examples:常见的B空间222222})，规定内积
\begin{align*}
(u,v)=\sum\limits_{|\alpha|\leqslant k}\int_\Omega \partial^\alpha u(x)\cdot\overline{\partial^\alpha v(x)}\mathrm{d}x\quad(\forall u,v\in C^k(\overline{\Omega})).
\end{align*}
那么 $(C^k(\overline{\Omega}),(\cdot,\cdot))$ 是一个内积空间.
\end{enumerate}
\end{proposition}
\begin{remark}
只有当$p=2$时,$L^p,l^p$中内积才满足\hyperref[proposition:平行四边形等式--泛函分析]{平行四边形法则},此时$L^p,l^p$才是Hilbert空间,其余情形均不是Hilbert空间.
\end{remark}
\begin{proof}
\end{proof}

\begin{lemma}[Poincaré不等式]\label{lemma:(庞加莱)Poincaré不等式}
设 $C_0^m(\Omega)$ 表示有界开区域 $\Omega\subset\mathbb{R}^n$ 上一切 $m$ 次连续可微, 并在边界 $\partial\Omega$ 的某邻域内为 $0$ 的函数集合, 即
\begin{align*}
C_0^m(\Omega)=\{u\in C^m(\overline{\Omega})\mid u(x)=0,\text{当}\ x\in\partial\Omega\ \text{的某邻域}\}.
\end{align*}
那么 $\forall u\in C_0^m(\Omega)$ 有
\begin{align}
\sum\limits_{|\alpha|<m}\int_\Omega |\partial^\alpha u(x)|^2\mathrm{d}x\leqslant C\sum\limits_{|\alpha|=m}\int_\Omega |\partial^\alpha u(x)|^2\mathrm{d}x,\label{eq:::--fiojoi3048HUAAIJOW9}
\end{align}
其中 $C$ 是仅依赖于区域 $\Omega$ 及 $m$ 的常数,$\alpha = (\alpha_1,\alpha_2,\cdots,\alpha_n), |\alpha| = \alpha_1+\alpha_2+\cdots+\alpha_n$，以及
\begin{align*}
\partial^\alpha u(x) = \frac{\partial^{|\alpha|}}{\partial x_1^{\alpha_1} \partial x_2^{\alpha_2} \cdots \partial x_n^{\alpha_n}} u(x).
\end{align*}
进而在 $C_0^m(\Omega)$ 上,定义
\begin{align}
\|u\|_m\triangleq\left(\sum\limits_{|\alpha|=m}\int_\Omega |\partial^\alpha u(x)|^2\mathrm{d}x\right)^\frac{1}{2}\label{eq:::--fiojoi3048HUAAIJOW12}
\end{align}
和
\begin{align*}
\|u\|\triangleq\left(\sum\limits_{|\alpha|\leqslant m}\int_\Omega |\partial^\alpha u(x)|^2\mathrm{d}x\right)^\frac{1}{2}
\end{align*}
是一对等价范数. 记 $C_0^m(\Omega)$ 按 \eqref{eq:::--fiojoi3048HUAAIJOW12} 式完备化后的空间为 $H_0^m(\Omega)$. 它是 $H^m(\Omega)$ 的一个闭子空间.
\end{lemma}
\begin{proof}
因为 $\Omega$ 是有界的, 我们可以把 $\Omega$ 放在某个边长为 $a$ 的立方体 $\Omega_1$ 内, 适当选择坐标系, 使得
\begin{align*}
\Omega_1=\{(x_1,x_2,\cdots,x_n)\in\mathbb{R}^n\mid 0\leqslant x_i\leqslant a\ (i=1,2,\cdots,n)\}.
\end{align*}
在 $\Omega_1\setminus\Omega$ 上补充定义 $u=0$, 经补充定义后, $u(x)$ 在 $\Omega_1$ 上 $m$ 次连续可微, 而且在边界上等于 $0$. $\forall x\in\Omega_1$,
\begin{align*}
u(x)=\int_0^{x_1}\frac{\partial u}{\partial t}(t,x_2,\cdots,x_n)\mathrm{d}t.
\end{align*}
再利用\hyperref[theorem:泛函分析--Cauchy-Schwarz不等式]{Cauchy-Schwarz不等式}, 我们有
\begin{align}
|u(x)|^2\leqslant \int_0^a{1^2\mathrm{d}x_1}\cdot \int_0^a{\left| \frac{\partial u}{\partial x_1} \right|^2\mathrm{d}x_1}=a\int_0^a{\left| \frac{\partial u}{\partial x_1} \right|^2\mathrm{d}x_1}.\label{eq:::--fiojoi3048HUAAIJOW10}
\end{align}
在 $\Omega_1$ 上对不等式 \eqref{eq:::--fiojoi3048HUAAIJOW10}积分得
\begin{align*}
\int_{\Omega}{\left| u\left( x \right) \right|^2\mathrm{d}x}&=\int_{\Omega _1}{\left| u\left( x \right) \right|^2\mathrm{d}x}\leqslant \int_{\Omega _1}{\left( a\int_0^a{\left| \frac{\partial u}{\partial x_1} \right|^2\mathrm{d}x_1} \right) \mathrm{d}x}
\\
&=\int_0^a{\cdots \int_0^a{\int_0^a{\left( a\int_0^a{\left| \frac{\partial u}{\partial x_1} \right|^2\mathrm{d}x_1} \right) \mathrm{d}x_1\mathrm{d}x_2\cdots \mathrm{d}x_n}}}
\\
&=\int_0^a{\cdots \int_0^a{\left( a^2\int_0^a{\left| \frac{\partial u}{\partial x_1} \right|^2\mathrm{d}x_1} \right) \mathrm{d}x_2\cdots \mathrm{d}x_n}}
\\
&=a^2\int_{\Omega _1}{\left| \frac{\partial u}{\partial x_1} \right|^2\mathrm{d}x}\leqslant a^2\int_{\Omega}{\left| \mathrm{grad}\,u\left( x \right) \right|^2\mathrm{d}x}
\\
&=a^2\int_{\Omega}{\sum_{i=1}^n{\left| \frac{\partial u}{\partial x_i} \right|^2\mathrm{d}x}}=a^2\int_{\Omega}{\sum_{\left| \alpha \right|=1}{\left| \partial ^{\alpha}u\left( x \right) \right|^2}\mathrm{d}x}
\\
&=a^2\sum_{\left| \alpha \right|=1}{\int_{\Omega}{\left| \partial ^{\alpha}u\left( x \right) \right|^2\mathrm{d}x}}.
\end{align*}
这就是\eqref{eq:::--fiojoi3048HUAAIJOW9}式中$m=1$的情形.然后逐次应用上式于 $\partial^\alpha u(x)\ (|\alpha|<m)$, 即得不等式 \eqref{eq:::--fiojoi3048HUAAIJOW9}.
\end{proof}

\begin{proposition}
$H_0^m(\Omega)$ 是一个 Hilbert 空间, 其内积定义为
\begin{align*}
(u,v)_m=\sum\limits_{|\alpha|=m}\int_\Omega \partial^\alpha u(x)\cdot\overline{\partial^\alpha v(x)}\mathrm{d}x\quad(\forall u,v\in C_0^m(\Omega)).
\end{align*}
\end{proposition}
\begin{remark}
当 $\Omega$ 的边界 $\partial\Omega$ 具有光滑的法向导数时, $H_0^m(\Omega)$ 中的元素实际上是满足边界条件:
\begin{align*}
u|_{\partial\Omega}=\frac{\partial u}{\partial n}\bigg|_{\partial\Omega}=\cdots=\left(\frac{\partial}{\partial n}\right)^{m-1}u\bigg|_{\partial\Omega}=0
\end{align*}
的 $C^m(\Omega)$ 函数 $u$ 的一种推广, 其中 $\frac{\partial}{\partial n}$ 是 $\partial\Omega$ 上的法向导数.
\end{remark}


\subsection{正交与正交基}

\begin{definition}
内积空间 $\mathscr{X}$ 上的两个元素 $x$ 与 $y$ 称为是\textbf{正交的}, 是指
\begin{align*}
(x,y)=0,
\end{align*}
记作 $x\perp y$. 又设 $M$ 是 $\mathscr{X}$ 的一个非空子集, $x\in\mathscr{X}$. 若对 $\forall y\in M$ 都有 $x\perp y$, 则称 $x$ 与 $M$ \textbf{正交}, 记作 $x\perp M$. 此外我们还称集合
\begin{align*}
\{x\in\mathscr{X}\mid x\perp M\}
\end{align*}
为 $M$ 的\textbf{正交补}, 记作 $M^\perp$.
\end{definition}

\begin{proposition}\label{proposition:正交的基本性质}
设 $\mathscr{X}$ 是内积空间, $M,N,S_1,\cdots,S_n$都是 $\mathscr{X}$的非空子集,
\begin{enumerate}[(1)]
\item\label{proposition:正交的基本性质-2} 若 $x=y+z$, 且 $y\perp z$, 则 $\|x\|^2=\|y\|^2+\|z\|^2$.

\item\label{proposition:正交的基本性质-3} 若 $x\perp y_n\ (n\in\mathbb{N})$, 且 $y_n\to y\,(n\to \infty)$, 则 $x\perp y$.

\item\label{proposition:正交的基本性质-4} $\bigcap_{k=1}^n{S_{k}^{\bot}}=\left( \bigcup_{k=1}^n{S_k} \right) ^{\bot},\quad \bigcup_{k=1}^n{S_{k}^{\bot}}\subset \left( \bigcap_{k=1}^n{S_k} \right) ^{\bot},\quad M^{\bot}=(\text{span}M)^{\bot}.$

\item\label{proposition:正交的基本性质-5} 若$M$是$\mathscr{X}$的线性子空间,则$M^\perp$ 是 $\mathscr{X}$ 的一个闭线性子空间,且$M\cap M^{\bot}=\{\theta\}$,进而
\begin{align*}
\mathscr{X}=M\oplus M^{\bot}.
\end{align*}
特别地,若$M$是$\mathscr{X}$的有限维线性子空间,则$M$一定是闭线性子空间.

\item\label{proposition:正交的基本性质-6} $M\subseteq N\implies N^\perp\subseteq M^\perp.$
\end{enumerate}
\end{proposition}
\begin{proof}
\end{proof}

\begin{proposition}\label{proposition:Hilbert空间稠密线性空间正交补为0}
设 $M$ 是 Hilbert 空间 $\mathscr{X}$ 的子集, 则
\begin{enumerate}[(1)]
\item\label{proposition:Hilbert空间稠密线性空间正交补为0-1} $(M^\perp)^\perp=\overline{\text{span}M}.$特别地,若$M$是有限维线性子空间,则$(M^\perp)^\perp=M$.

\item\label{proposition:Hilbert空间稠密线性空间正交补为0-2} 若$M$为$\mathscr{X}$的线性子空间，则
$M^\perp = \{\theta\}$ 当且仅当 $\overline{M} = \mathscr{X}$（即 $M$ 在 $\mathscr{X}$ 中稠密）。
\end{enumerate}
\end{proposition}
\begin{proof}
\begin{enumerate}[(1)]
\item 先证 $\overline{\operatorname{span}M}\subseteq (M^\perp)^\perp$. 
任取 $x\in M$，对任意 $y\in M^\perp$，有 $\langle x,y\rangle=0$，故 $x\perp M^\perp$，即 $x\in (M^\perp)^\perp$. 
因此 $M\subseteq (M^\perp)^\perp$. 
由于 $(M^\perp)^\perp$ 是闭线性子空间，它包含 $\operatorname{span}M$,也包含其闭包，即 $\overline{\operatorname{span}M}\subseteq (M^\perp)^\perp$.

再证 $(M^\perp)^\perp\subseteq \overline{\operatorname{span}M}$. 
令 $V=\overline{\operatorname{span}M}$，则$V$ 是闭线性子空间，且由\rrefpro{proposition:正交的基本性质}{proposition:正交的基本性质-5}知有正交分解 $\mathscr{X}=V\oplus V^\perp$. 
任取 $x\in (M^\perp)^\perp$，则存在唯一分解 $x=v+w$，其中 $v\in V$，$w\in V^\perp$. 
由于 $M\subseteq V$，故由\rrefpro{proposition:正交的基本性质}{proposition:正交的基本性质-6}知$V^\perp\subseteq M^\perp$，于是 $w\in M^\perp$. 
又因 $x\in (M^\perp)^\perp$ 且 $w\in M^\perp$，有 $\langle x,w\rangle=0$. 
但
\begin{align*}
\langle x,w\rangle=\langle v+w,w\rangle=\langle v,w\rangle+\langle w,w\rangle=0+\|w\|^2=\|w\|^2\Longrightarrow \|w\|^2=0\Longrightarrow w=0.
\end{align*}
因此 $x=v\in V=\overline{\operatorname{span}M}$. 

特别地,由\rrefpro{proposition:正交的基本性质}{proposition:正交的基本性质-5}知$M$是闭线性子空间,从而$\overline{\operatorname{span}M}=M$.

\item 若 $\overline{M} = \mathscr{X}$，任取 $z \in M^\perp$，则对$\forall x \in \mathscr{X}$，存在序列 $\{m_n\} \subset M$ 使得 $m_n \to x$。由于 $(m_n, z)=0$，由\hyperref[proposition:内积的连续性]{内积的连续性}得 $(x, z)=0$，故 $z=\theta$，从而 $M^\perp = \{\theta\}$.

反过来，若 $\overline{M} \ne \mathscr{X}$，则 $\overline{M}$ 是 $\mathscr{X}$ 的闭真子空间。由正交补空间的定义知$\mathscr{X} = \overline{M} \oplus (\overline{M})^\perp$，且因为 $\overline{M} \ne \mathscr{X}$，有 $(\overline{M})^\perp \ne \{\theta\}$。取非零 $z \in (\overline{M})^\perp$,又$M\subset \overline{M}$,故 $z \in M^\perp$，因此 $M^\perp \ne \{\theta\}$,矛盾!
\end{enumerate}
\end{proof}

\begin{definition}
设 $\mathscr{X}$ 是一个内积空间,$\theta$是零向量, 集合 $S=\{e_\alpha\mid\alpha\in A\}$ 是 $\mathscr{X}$ 的一个子集. 称 $S$ 为\textbf{正交集}, 是指
\begin{align*}
e_\alpha\perp e_\beta\quad(\text{当}\ \alpha\neq\beta,\forall\alpha,\beta\in A).
\end{align*}
如果还有 $\|e_\alpha\|=(e_\alpha,e_\alpha)^{\frac{1}{2}}=1\ (\forall\alpha\in A)$, 则称 $S$ 为\textbf{正交规范集}. 又如果在 $\mathscr{X}$ 中不存在非零元与 $S$ 正交, 即 $S^\perp=\{\theta\}$, 那么称 $S$ 为\textbf{完备的}.
\end{definition}

\begin{lemma}[Zorn引理]\label{lemma:泛函分析--Zorn引理}
设 $\mathscr{X}$ 是一个半序集. 如果它的每一个全序子集有一个上界, 那么 $\mathscr{X}$ 有一个极大元.
\end{lemma}

\begin{proposition}
非 $\{\theta\}$ 内积空间 $\mathscr{X}$ 中必存在完备正交集.
\end{proposition}
\begin{proof}
因为 $\mathscr{X}\neq\{\theta\}$, 所以 $\mathscr{X}$ 中的正交集依包含关系构成一个半序集类$\mathcal{F} =\left\{ A\subseteq \mathscr{X} \mid A\text{是正交集} \right\} $, 并且$\mathcal{F}$的每个全序子集类$\mathcal{T} \subseteq \mathcal{F}$有一个上界, 就是这些集之并集$\bigcup_{A\subseteq \mathcal{T}}{A}$. 依\hyperref[lemma:泛函分析--Zorn引理]{Zorn 引理}, 这个半序集类有极大元. 我们来证明: 这个极大元 (记作 $S$) 就是完备正交集. 因若不然, 则必 $\exists x_0\in S^\perp,x_0\neq\theta$, 令 $S_1=\{x_0\}\cup S$, 得到 $S_1$ 还是正交集, 并且 $S\subsetneq S_1$, 这便与 $S$ 的极大性相矛盾!
\end{proof}

\begin{definition}
内积空间 $\mathscr{X}$ 中的正交规范集 $S=\{e_\alpha\mid\alpha\in A\}$, 称为一个\textbf{(正交)基} 或\textbf{封闭的} 是指 $\forall x\in\mathscr{X}$, 有下列表示:
\begin{align*}
x=\sum\limits_{\alpha\in A}(x,e_\alpha)e_\alpha,
\end{align*}
其中 $\{(x,e_\alpha)\mid\alpha\in A\}$ 称为 $x$ 关于基 $\{e_\alpha\mid\alpha\in A\}$ 的 $\mathbf{Fourier}$ \textbf{系数}.
\end{definition}

\begin{theorem}[Bessel不等式]\label{theorem:泛函分析---Bessel不等式}
设 $\mathscr{X}$ 是一个内积空间. 如果 $S=\{e_\alpha\mid\alpha\in A\}$ 是 $\mathscr{X}$ 中的正交规范集, 那么 $\forall x\in\mathscr{X}$, 有
\begin{align}
\sum\limits_{\alpha\in A}|(x,e_\alpha)|^2\leqslant\|x\|^2.\label{eq:::--fiojoi3048HUAAIJOW16}
\end{align}
并且满足$(x,e_\alpha)\neq 0$ 的 $\alpha\in A$至多有可数多个. 
\end{theorem}
\begin{proof}
首先对 $A$ 的任意有限子集, 不妨设它们是 $1,2,\cdots,n$, 证明
\begin{align}
\sum\limits_{i=1}^n|(x,e_i)|^2\leqslant\|x\|^2.\label{eq:::--fiojoi3048HUAAIJOW17}
\end{align}
因为
\begin{align*}
0\leqslant\left\|x-\sum\limits_{i=1}^n(x,e_i)e_i\right\|^2
=\left(x-\sum\limits_{i=1}^n(x,e_i)e_i,x-\sum\limits_{j=1}^n(x,e_j)e_j\right)=\|x\|^2-\sum\limits_{i=1}^n|(x,e_i)|^2,
\end{align*}
所以 \eqref{eq:::--fiojoi3048HUAAIJOW17} 式成立. 由此可见, 对 $\forall n\in\mathbb{N}$, 适合 $|(x,e_\alpha)|>\frac{1}{n}$ 的 $\alpha\in A$ 至多只有有限多个.否则有无穷多个$\alpha \in A$满足$\left| \left( x,e_{\alpha} \right) \right|>\frac{1}{n}$，从而存在$\left\{ \alpha _k \right\} _{k=1}^{\infty}\subseteq A$，使得
\begin{align*}
\left| \left( x,e_{\alpha _k} \right) \right|>\frac{1}{n},\quad \forall k\in \mathbb{N} \Longrightarrow \sum\limits_{k=1}^{\infty}{\left| \left( x,e_{\alpha _k} \right) \right|^2}\geqslant \sum\limits_{k=1}^{\infty}{\frac{1}{n^2}}=+\infty >\left\| x \right\| ^2,
\end{align*}
矛盾!因此
\begin{align*}
A_n=\{\alpha \in A\mid \left| \left( x,e_{\alpha} \right) \right|>\frac{1}{n}\},\quad \forall n\in \mathbb{N}
\end{align*}
都是有限集.从而
\begin{align*}
\bigcup_{n=1}^{\infty}{A_n}=\left\{ \alpha \in A\mid \left| \left( x,e_{\alpha} \right) \right|>0 \right\} 
\end{align*}
是至多可数集,
即满足$(x,e_\alpha)\neq 0$ 的 $\alpha\in A$至多有可数多个. 从而可设$\bigcup_{n=1}^{\infty}{A_n}=\{\alpha_i\}_{i=1}^{\infty}$.又由 \eqref{eq:::--fiojoi3048HUAAIJOW17} 式我们有
\begin{align*}
\sum_{i=1}^n{\left| \left( x,e_{\alpha _i} \right) \right|^2}\leqslant \left\| x \right\| ^2,\quad \forall n\in \mathbb{N} .
\end{align*}
于是
\begin{align*}
\sum_{\alpha \in A}{\left| \left( x,e_{\alpha} \right) \right|^2}=\sum_{\alpha \in \bigcup_{n=1}^{\infty}{A_n}}{\left| \left( x,e_{\alpha} \right) \right|^2}=\sum_{i=1}^{\infty}{\left| \left( x,e_{\alpha _i} \right) \right|^2}=\underset{n\in \mathbb{N}}{\mathrm{sup}}\sum_{i=1}^n{\left| \left( x,e_{\alpha _i} \right) \right|^2}\leqslant \left\| x \right\| ^2.
\end{align*}
\end{proof}

\begin{corollary}\label{corollary:正交规范集恒等式}
假设 $\mathscr{X}$ 是 Hilbert 空间, 且 $\{e_\alpha\mid\alpha\in A\}$ 是 $\mathscr{X}$ 中的正交规范集. 那么对 $\forall x\in\mathscr{X}$, 有
\begin{align*}
\sum\limits_{\alpha\in A}(x,e_\alpha)e_\alpha\in\mathscr{X},
\end{align*}
且
\begin{align}
\left\|x-\sum\limits_{\alpha\in A}(x,e_\alpha)e_\alpha\right\|^2=\|x\|^2-\sum\limits_{\alpha\in A}|(x,e_\alpha)|^2.\label{eq:::--fiojoi3048HUAAIJOW18}
\end{align}
\end{corollary}
\begin{proof}
由\hyperref[theorem:泛函分析---Bessel不等式]{Bessel不等式}可不妨设使得 $(x,e_\alpha)\neq 0$ 的可数多个 $\alpha\in A$ 是 $1,2,\cdots,n,\cdots$, 那么
\begin{align*}
\sum\limits_{\alpha\in A}(x,e_\alpha)e_\alpha=\sum\limits_{n=1}^\infty(x,e_n)e_n,
\end{align*}
并由\hyperref[theorem:泛函分析---Bessel不等式]{Bessel不等式}可知 $\sum\limits_{n=1}^\infty|(x,e_n)|^2$ 收敛, 因此有
\begin{align*}
\left\|\sum\limits_{n=m}^{m+p}(x,e_n)e_n\right\|^2=\sum\limits_{n=m}^{m+p}|(x,e_n)|^2\to 0\quad(m\to\infty,\forall p\in\mathbb{N}).
\end{align*}
于是, $\left\{x_m=\sum\limits_{n=1}^m(x,e_n)e_n\right\}$ 是基本列, 从而
\begin{align*}
\sum\limits_{\alpha\in A}(x,e_\alpha)e_\alpha=\sum\limits_{n=1}^\infty(x,e_n)e_n=\lim_{m\to\infty}x_m\in\mathscr{X}.
\end{align*}
又因为对$\forall m\in \mathbb{N}$,有
\begin{align*}
\left( x-\sum_{n=1}^m{\left( x,e_n \right) e_n},\sum_{n=1}^m{\left( x,e_n \right) e_n} \right) &=\left( x,\sum_{n=1}^m{\left( x,e_n \right) e_n} \right) -\left( \sum_{n=1}^m{\left( x,e_n \right) e_n},\sum_{n=1}^m{\left( x,e_n \right) e_n} \right) 
\\
&=\sum_{n=1}^m{\left( x,\left( x,e_n \right) e_n \right)}-\sum_{n=1}^m{\left( \left( x,e_n \right) e_n,\sum_{n=1}^m{\left( x,e_n \right) e_n} \right)}
\\
&=\sum_{n=1}^m{\overline{\left( x,e_n \right) }\left( x,e_n \right)}-\sum_{n=1}^m{\left[ \left( x,e_n \right) \sum_{k=1}^m{\left( e_n,\left( x,e_k \right) e_k \right)} \right]}
\\
&=\sum_{n=1}^m{\left| \left( x,e_n \right) \right|^2}-\sum_{n=1}^m{\left[ \left( x,e_n \right) \overline{\left( x,e_n \right) }\left( e_n,e_n \right) \right]}
\\
&=\sum_{n=1}^m{\left| \left( x,e_n \right) \right|^2}-\sum_{n=1}^m{\left| \left( x,e_n \right) \right|^2}=0,
\end{align*}
即
\begin{align*}
\left( x-\sum_{n=1}^m{(x,e_n)e_n} \right) \bot \sum_{n=1}^m{(x,e_n)e_n}.
\end{align*}
由\rrefpro{proposition:正交的基本性质}{proposition:正交的基本性质-3}知$\left( x-\sum_{n=1}^{\infty}{(x,e_n)e_n} \right) \bot \sum_{n=1}^{\infty}{(x,e_n)e_n}$.
所以再\rrefpro{proposition:正交的基本性质}{proposition:正交的基本性质-2}知
\begin{align*}
\left\|x-\sum\limits_{n=1}^\infty(x,e_n)e_n\right\|^2+\sum\limits_{n=1}^\infty|(x,e_n)|^2=\|x\|^2.
\end{align*}
这就是 \eqref{eq:::--fiojoi3048HUAAIJOW18} 式.
\end{proof}

\begin{theorem}[Parseval恒等式]\label{theorem:泛函分析--Parseval恒等式}
设$\mathscr{X}$是一个Hilbert空间，若$S=\{e_\alpha \mid \alpha \in A\}$是$\mathscr{X}$中的正交规范集，则如下三条等价：
\begin{enumerate}[(1)]
\item $S$是封闭的；

\item $S$是完备的；

\item $\mathbf{Parseval}$\textbf{恒等式}
\begin{align}
\|x\|^2 = \sum\limits_{\alpha \in A} |(x, e_\alpha)|^2 \quad (\forall x \in \mathscr{X}). \label{eq:parseval}
\end{align}
成立。
\end{enumerate}
\end{theorem}
\begin{proof}

(1)$\implies$(2). 若$S$不完备，则$\exists x \in \mathscr{X} \setminus \{\theta\}$，使得
\begin{align*}
(x, e_\alpha) = 0 \quad (\forall \alpha \in A).
\end{align*}
但由封闭性有$x = \sum\limits_{\alpha \in A} (x, e_\alpha) e_\alpha = \theta$，矛盾。

(2)$\implies$(3). 若$\exists x \in \mathscr{X}$使Parseval等式\eqref{eq:parseval}不成立，则由\refcor{corollary:正交规范集恒等式}知
\begin{align*}
\left\|x - \sum\limits_{\alpha \in A} (x, e_\alpha) e_\alpha\right\|^2 = \|x\|^2 - \sum\limits_{\alpha \in A} |(x, e_\alpha)|^2 > 0.
\end{align*}
于是$y \triangleq x - \sum\limits_{\alpha \in A} (x, e_\alpha) e_\alpha \neq \theta$.又注意到
\begin{align*}
\left( y,e_{\beta} \right) =\left( x-\sum_{\alpha \in A}{\left( x,e_{\alpha} \right) e_{\alpha}},e_{\beta} \right) =\left( x,e_{\beta} \right) -\left( \sum_{\alpha \in A}{\left( x,e_{\alpha} \right) e_{\alpha}},e_{\beta} \right) =\left( x,e_{\beta} \right) -\left( x,e_{\beta} \right) \left( e_{\beta},e_{\beta} \right) =0,\quad \forall \beta \in A.
\end{align*}
故$y \in S^\perp$。这与$S$完备性矛盾。

(3)$\implies$(1). 对$\forall x\in \mathscr{X}$,联合Parseval等式\eqref{eq:parseval}与\refcor{corollary:正交规范集恒等式}知
\begin{align*}
\left\|x - \sum\limits_{\alpha \in A} (x, e_\alpha) e_\alpha\right\|^2 = \|x\|^2 - \sum\limits_{\alpha \in A} |(x, e_\alpha)|^2 = 0.
\end{align*}
因此有
\begin{align*}
x = \sum\limits_{\alpha \in A} (x, e_\alpha) e_\alpha.
\end{align*}
故$S$是封闭的.
\end{proof}

\begin{proposition}\label{proposition:H2D空间--泛函分析}
\begin{enumerate}[(1)]
\item\label{proposition:H2D空间--泛函分析-1} 在$L^2[0, 2\pi]$上，
\begin{align*}
e_n(t) = \frac{1}{\sqrt{2\pi}} e^{\mathrm{i}nt} \quad (n=0, \pm1, \pm2, \cdots)
\end{align*}
是一组正交规范基。$\forall u \in L^2[0, 2\pi]$，对应的Fourier系数是
\begin{align*}
(u, e_n) = \frac{1}{\sqrt{2\pi}} \int_{0}^{2\pi} u(t) e^{-\mathrm{i}nt} \mathrm{d}t \quad (n=0, \pm1, \pm2, \cdots).
\end{align*}

\item\label{proposition:H2D空间--泛函分析-2} 在$l^2$空间上，
\begin{align*}
e_n = (\underbrace{0, 0, \cdots, 0, 1}_{n}, 0, \cdots) \quad (n=1,2,3,\cdots)
\end{align*}
是一组正交规范基。

\item\label{proposition:H2D空间--泛函分析-3} 设$D$是$\mathbb{C}$中的单位开圆域，$H^2(D)$表示在$D$内满足
\begin{align*}
\iint_{D} |u(z)|^2 \mathrm{d}x\mathrm{d}y < \infty \quad (z=x+\mathrm{i}y)
\end{align*}
的解析函数全体组成的空间。规定内积为
\begin{align*}
(u, v) = \iint_{D} u(z) \overline{v(z)} \mathrm{d}x\mathrm{d}y.
\end{align*}
这时函数组
\begin{align*}
\varphi_n(z) = \sqrt{\frac{n}{\pi}} z^{n-1} \quad (n=1,2,3,\cdots)
\end{align*}
是一组正交规范基。并且若$u(z) = \sum\limits_{k=0}^{\infty} b_k z^k \in H^2(D)$，则它对应的Fourier系数是
\begin{align*}
(u, \varphi_n)  = b_{n-1} \sqrt{\frac{\pi}{n}} \quad (n=1,2,3,\cdots).
\end{align*}
\end{enumerate}
\end{proposition}
\begin{enumerate}[(1)]
\item 

\item 

\item 先证明 $\{\varphi_n\}$ 是标准正交系。任取正整数 $m,n$，利用极坐标 $z=r e^{\mathrm{i}\theta}$，$\mathrm{d}x\mathrm{d}y=r\,\mathrm{d}r\mathrm{d}\theta$，有
\begin{align}
(\varphi_m,\varphi_n)
&=\iint_D \sqrt{\frac{m}{\pi}}z^{m-1}\overline{\sqrt{\frac{n}{\pi}}z^{n-1}}\,\mathrm{d}x\mathrm{d}y \nonumber \\
&=\frac{\sqrt{mn}}{\pi}\int_0^1\int_0^{2\pi} r^{(m-1)+(n-1)} e^{\mathrm{i}(m-n)\theta}\, r\,\mathrm{d}\theta\,\mathrm{d}r \nonumber \\
&=\frac{\sqrt{mn}}{\pi}\int_0^1 r^{m+n-1}\,\mathrm{d}r \int_0^{2\pi} e^{\mathrm{i}(m-n)\theta}\,\mathrm{d}\theta.  \label{eq::iofjiwjo39023}
\end{align}
当 $m\neq n$ 时，$\int_0^{2\pi} e^{\mathrm{i}(m-n)\theta}\,\mathrm{d}\theta=0$，故 $(\varphi_m,\varphi_n)=0$. 当 $m=n$ 时，
\begin{align}
\|\varphi_n\|^2
=(\varphi_n,\varphi_n)=\frac{n}{\pi}\int_0^1 r^{2n-1}\,\mathrm{d}r\int_0^{2\pi}\mathrm{d}\theta
=\frac{n}{\pi}\cdot \frac{1}{2n}\cdot 2\pi =1. \label{eq::iofjiwjo39023-1}
\end{align}
因此 $\{\varphi_n\}_{n=1}^{\infty}$ 是标准正交系。

下面证明完备性。任取 $u\in H^2(D)$，由于 $u$ 在单位圆盘内解析，故由\hyperref[Complex Analysis-theorem:全纯函数的Taylor展开]{全纯函数的Taylor展开}知,$u$可展开为 Taylor 级数
\begin{align*}
u(z)=\sum_{k=0}^{\infty} b_k z^k,\quad |z|<1. 
\end{align*}
对任意 $n\geqslant 1$，由\hyperref[proposition:内积的连续性]{内积的连续性}得
\begin{align*}
(u,\varphi _n)&=\left( \sum_{k=0}^{\infty}{b_kz^k},\varphi _n \right) =\left( \underset{N\rightarrow \infty}{\lim}\sum_{k=0}^N{b_kz^k},\varphi _n \right)\\
&=\underset{N\rightarrow \infty}{\lim}\left( \sum_{k=0}^N{b_kz^k},\varphi _n \right) =\underset{N\rightarrow \infty}{\lim}\iint_D{\left( \sum_{k=0}^N{b_kz^k} \right) \overline{\sqrt{\frac{n}{\pi}}z^{n-1}}\,\mathrm{d}x\mathrm{d}y}\\
&=\underset{N\rightarrow \infty}{\lim}\sqrt{\frac{n}{\pi}}\sum_{k=0}^N{b_k\iint_D{z^k\overline{z^{\,n-1}}\,\mathrm{d}x\mathrm{d}y}}=\sqrt{\frac{n}{\pi}}\sum_{k=0}^{\infty}{b_k\iint_D{z^k\overline{z^{\,n-1}}\,\mathrm{d}x\mathrm{d}y}}\\
&\xlongequal{\text{\eqref{eq::iofjiwjo39023}\eqref{eq::iofjiwjo39023-1}式}}\sqrt{\frac{n}{\pi}}\sum_{k=0}^{\infty}{b_k}\cdot \frac{\pi}{k+1}\delta _{k,n-1}=\sqrt{\frac{n}{\pi}}\cdot b_{n-1}\cdot \frac{\pi}{n}=\sqrt{\frac{\pi}{n}}b_{n-1},
\end{align*}
其中$\delta_{k,n-1}$是Kronecker符号.若 $u$ 与所有 $\varphi_n$ 都正交,则$(u,\varphi _n)=0$.由上式知$b_{n-1}=0$ 对所有 $n$ 成立，即所有 $b_k=0$，故 $u\equiv 0$. 这表明 $\{\varphi_n\}$ 的正交补空间只有零向量，所以它是完备的，从而是 $H^2(D)$ 的一组标准正交基。
\end{enumerate}


\subsection{正交化和Hilbert空间的同构}

\begin{theorem}[Gram-Schmidt 正交化]\label{theorem:Gram-Schmidt正交化}
设 $\mathscr{X}$ 是内积空间，$\{x_n\}_{n=1}^{\infty} \subset \mathscr{X}$ 是线性无关的向量列。则存在正交规范列 $\{y_n\}_{n=1}^{\infty} \subset \mathscr{X}$，使得对任意 $n \in \mathbb{N}$ 有
\[
\operatorname{span}\{y_1,\cdots,y_n\} = \operatorname{span}\{x_1,\cdots,x_n\}.
\]
且满足
\[
\begin{cases}
y_1 = x_1, \quad e_1 = \dfrac{y_1}{\|y_1\|},\\[6pt]
y_n = x_n - \displaystyle\sum_{i=1}^{n-1} \langle x_n, e_i \rangle e_i, \quad e_n = \dfrac{y_n}{\|y_n\|}, \quad n \geqslant 2.
\end{cases}
\]
这里 $\langle\cdot,\cdot\rangle$ 与 $\|\cdot\|$ 分别为 $\mathscr{X}$ 上的内积和内积诱导的范数。

若 $\{x_n\}$ 张成 $\mathscr{X}$ 的某个子空间 $M$，则 $\{y_n\}$ 构成 $M$ 的一组标准正交基。
\end{theorem}

\begin{definition}
设$(\mathscr{X}_1, (\cdot, \cdot)_1)$与$(\mathscr{X}_2, (\cdot, \cdot)_2)$是两个内积空间，如果存在$\mathscr{X}_1 \to \mathscr{X}_2$的一个线性同构$T$，满足
\begin{align*}
(Tx, Ty)_2 = (x, y)_1 \quad (\forall x, y \in \mathscr{X}_1),
\end{align*}
则称内积空间$\mathscr{X}_1$与$\mathscr{X}_2$是\textbf{同构的}。
\end{definition}

\begin{theorem}
Hilbert空间$\mathscr{X}$是可分的当且仅当它有至多可数的正交规范基$S$。

当$S$是$\mathscr{X}$的正交规范基时,若$S$的元素个数$N < \infty$，则$\mathscr{X}$同构于$\mathbb{K}^N(\mathbb{K}=\mathbb{R}\text{或}\mathbb{C})$;
若$S$的元素个数$N = \infty$，则$\mathscr{X}$同构于$l^2$。
\end{theorem}
\begin{proof}

{\heiti 必要性:} 设$\{x_n\}_{n=1}^{\infty}$是$\mathscr{X}$中的可数稠密子集，令$y_1=x_1$，若$\mathrm{span}\{ y_1 \} = \mathrm{span}\{ x_n \}_{n=1}^{\infty}$，则$\{ y_1 \}$就是$\{ x_n \}_{n=1}^{\infty}$的线性无关子集。
否则$\mathrm{span}\{ y_1 \} \subset \mathrm{span}\{ x_n \}_{n=1}^{\infty}$，则任取$y_2\in \mathrm{span}\{ x_n \}_{n=1}^{\infty}\setminus \mathrm{span}\{ y_1 \}$，若$\mathrm{span}\{ y_1,y_2 \} = \mathrm{span}\{ x_n \}_{n=1}^{\infty}$，则$\{ y_1,y_2 \}$就是$\{ x_n \}_{n=1}^{\infty}$的线性无关子集；
否则$\mathrm{span}\{ y_1,y_2 \} \subset \mathrm{span}\{ x_n \}_{n=1}^{\infty}$，则任取$y_3\in \mathrm{span}\{ x_n \}_{n=1}^{\infty}\setminus \mathrm{span}\{ y_1,y_2 \}$，再考虑$\mathrm{span}\{ y_1,y_2,y_3 \}$。
依此类推，最终得到$\{ x_n \}_{n=1}^{\infty}$的线性无关子集$\{ y_n \}_{n=1}^{N}$（$N<\infty$或$N=\infty$），满足
\begin{align*}
\overline{\mathrm{span}\{y_n\}_{n=1}^{N}}=\overline{\mathrm{span}\{x_n\}_{n=1}^{\infty}}=\mathscr{X}.
\end{align*}
再对$\{y_n\}_{n=1}^{N}$应用\hyperref[theorem:Gram-Schmidt正交化]{Gram-Schmidt正交化}，便构造出一个正交规范集$\{e_n\}_{n=1}^{N}$。又因为
\begin{align*}
\overline{\operatorname{span}\{e_n\}_{n=1}^{N}} = \overline{\operatorname{span}\{y_n\}_{n=1}^{N}} = \mathscr{X},
\end{align*}
所以$\{e_n\}_{n=1}^{N}$是正交规范基。

{\heiti 充分性:} 设$\{e_n\}_{n=1}^{N}$（$N < \infty$或$N = \infty$）是$\mathscr{X}$的正交规范基，那么集合
\begin{align*}
S\triangleq \left\{x = \sum\limits_{n=1}^{N} a_n e_n \mid \operatorname{Re}a_n \text{ 与 } \operatorname{Im}a_n \text{ 皆为有理数}\right\}
\end{align*}
因为$\mathbb{Q}$和$\{ e_n \}_{n=1}^{N}$都是至多可数集，所以集合$S$也是可数集。
对$\forall x\in \mathscr{X}$，由$\{ e_n \}_{n=1}^{N}$是$\mathscr{X}$的正交规范基知
\begin{align*}
x=\sum\limits_{n=1}^N (x,e_n) e_n.
\end{align*}
不妨设$(x,e_n)=x_n+\mathrm{i}y_n\in \mathbb{C}$，其中$x_n,y_n\in \mathbb{R}$，从而存在$\{ p_{n}^{(k)} \}_{k=1}^{\infty},\{ q_{n}^{(k)} \}_{k=1}^{\infty}\subseteq \mathbb{Q}$，使得
\begin{align*}
p_{n}^{(k)}\rightarrow x_n,\quad q_{n}^{(k)}\rightarrow y_n \quad (k\rightarrow \infty).
\end{align*}
令$r_{n}^{(k)}=p_{n}^{(k)}+\mathrm{i}q_{n}^{(k)},\forall k,n\in \mathbb{N}$，则$r_{n}^{(k)}\rightarrow (x,e_n) \ (k\rightarrow \infty)$。再令$r^{(k)}=\sum\limits_{n=1}^N r_{n}^{(k)} e_n,\forall k\in \mathbb{N}$，则
\begin{align*}
r^{(k)}\in S \ (\forall k\in \mathbb{N}) \quad \text{且} \quad r^{(k)}\rightarrow \sum\limits_{n=1}^N (x,e_n) e_n=x \ (k\rightarrow \infty).
\end{align*}
因此$S$是$\mathscr{X}$的可数稠密子集。从而$\mathscr{X}$是可分的。

当$S$是$\mathscr{X}$的正交规范基时,对于正交规范基$\{e_n\}_{n=1}^{N}$（$N < \infty$或$N = \infty$），做对应
\begin{align*}
T: x \mapsto \{(x, e_n)\}_{n=1}^{N} \quad (\forall x \in \mathscr{X}).
\end{align*}
显然$T$是$\mathscr{X} \to \mathbb{K}^N$（当$N < \infty$）或$\mathscr{X} \to l^2$（当$N = \infty$）的线性映射.记
\begin{align*}
\xi_n = (\underbrace{0, 0, \cdots, 0, 1}_{n}, 0, \cdots) \quad (n=1,2,3,\cdots).
\end{align*}
则当$N<\infty$时,$\{\xi_n\}_{n=1}^{N}$是$\mathbb{K}^N$的正交规范基;当$N=\infty$时,由\rrefpro{proposition:H2D空间--泛函分析}{proposition:H2D空间--泛函分析-2}知$\{\xi_n\}_{n=1}^{\infty}$是$l^2$的正交规范基.而无论$N<\infty$或$N=\infty$,都有
\begin{align*}
T(e_n) = \xi_n\quad (n=1,2,3,\cdots).
\end{align*}
故$T$是双射.
此外,
\begin{align*}
(x, y) = \left( \sum\limits_{i=1}^{N} (x, e_i) e_i, \sum\limits_{j=1}^{N} (y, e_j) e_j \right) 
= \sum\limits_{i=1}^{N} (x, e_i) \overline{(y, e_i)} \quad (\forall x, y \in \mathscr{X}).
\end{align*}
因此$T$还是保持内积的。于是当$N < \infty$时，$\mathscr{X}$同构于$\mathbb{K}^N$；而当$N = \infty$时，$\mathscr{X}$同构于$l^2$。
\end{proof}


\subsection{再论最佳逼近问题}

\begin{theorem}\label{theorem:泛函分析---定理1.1611}
如果$C$是Hilbert空间$\mathscr{X}$中的一个闭凸子集，那么在$C$上存在唯一元素$x_0$取到最小范数,即$\|x_0\|=\inf\limits_{z \in C} \|z\|.$
\end{theorem}
\begin{proof}

{\heiti 存在性:} 若零向量$\theta \in C$，则$x_0 = \theta$；若零向量$\theta \notin C$，则
\begin{align*}
d \triangleq \inf\limits_{z \in C} \|z\| > 0.
\end{align*}
由下确界定义，$\forall n \in \mathbb{N}$，$\exists x_n \in C$，使得
\begin{align}
d \leqslant \|x_n\| \leqslant d + \frac{1}{n} \quad (n=1,2,\cdots). \label{eq:min_norm_seq}
\end{align}
对$\forall m,n\in \mathbb{N}$且$m>n$,利用\hyperref[proposition:平行四边形等式--泛函分析]{平行四边形等式}得
\begin{align*}
\|x_m - x_n\|^2 = 2(\|x_m\|^2 + \|x_n\|^2) - 4\left\|\frac{x_m + x_n}{2}\right\|^2 
\leqslant 2\left[ \left(d + \frac{1}{n}\right)^2 + \left(d + \frac{1}{m}\right)^2 \right] - 4d^2 \to 0 \quad (\text{当 } n, m \to \infty).
\end{align*}
故$\{x_n\}$是一个基本列,从而$\{x_n\}$有极限$x_0$.由$C$的闭性，$x_0 \in C$，并由\eqref{eq:min_norm_seq}式，$\|x_0\| = d$，即$x_0$取到了$C$上元素的最小范数。

{\heiti 唯一性:} 如果有$x_0, \hat{x}_0 \in C$使得$\|x_0\| = \|\hat{x}_0\| = d=\underset{z\in C}{\min}\left\| z \right\| $，由$C$是凸集知$\frac{x_0+\hat{x}_0}{2}\in C$,那么
\begin{align*}
\|x_0 - \hat{x}_0\|^2 = 2(\|x_0\|^2 + \|\hat{x}_0\|^2) - 4\left\|\frac{x_0 + \hat{x}_0}{2}\right\|^2 
\leqslant 4d^2 - 4d^2 = 0,
\end{align*}
即得$x_0 = \hat{x}_0$。

\begin{figure}[H]
\centering
\includegraphics[scale=0.3]{泛函分析图1.6.1.png}
\caption{}
\label{figure:min_norm}
\end{figure}
\end{proof}

\begin{corollary}[最佳逼近元的存在唯一性]\label{corollary:泛函分析--最佳逼近元的存在唯一性}
若$C$是Hilbert空间$\mathscr{X}$中的一个闭凸子集，则对$\forall y \in \mathscr{X}$，都存在唯一的最佳逼近元$x_0 \in C$，使得
\begin{align}
\|y - x_0\| = \inf\limits_{x \in C} \|x - y\|. \label{eq:best_approx}
\end{align}
特别地,若$M$是Hilbert空间$\mathscr{X}$上的一个闭线性子空间，则对$\forall y \in \mathscr{X}$，都存在唯一的最佳逼近元$ x_0 \in M$，使得
\begin{align*}
\|y - x_0\| = \inf\limits_{x \in M} \|x - y\|.
\end{align*}
\end{corollary}
\begin{proof}
只需考察集合$C - \{y\} \triangleq \{x - y \mid x \in C\}$，它显然还是$\mathscr{X}$上的一个闭凸子集。根据\refthe{theorem:泛函分析---定理1.1611}，存在唯一的$ z_0 \in C - \{y\}$，使得
\begin{align*}
\parallel z_0\parallel =\mathop {\mathrm{inf}} \limits_{z\in C-\{y\}}\parallel z\parallel =\mathop {\mathrm{inf}} \limits_{z\in C}\parallel z-y\parallel .
\end{align*}
令$x_0 \triangleq z_0 + y$，便得到\eqref{eq:best_approx}式。

特别地,显然闭线性子空间$M$是$\mathscr{X}$是上的一个闭凸子集,故由上述证明知结论成立.
\end{proof}

\begin{theorem}[最佳逼近元]\label{theorem:泛函分析--最佳逼近元的刻画}
设$C$是内积空间$\mathscr{X}$中的一个闭凸子集，$\forall y \in \mathscr{X}$，则$x_0$是$y$在$C$上的最佳逼近元当且仅当它适合
\begin{align}
\operatorname{Re}(y - x_0, x_0 - x) \geqslant 0 \quad (\forall x \in C). \label{eq:best_approx_char}
\end{align}
\end{theorem}
\begin{proof}
对$\forall x \in C$，考察函数
\begin{align*}
\varphi_x(t) = \|y - tx - (1-t)x_0\|^2 \quad (t \in [0,1]).
\end{align*}
事实上,若$x_0$是$y$在$C$上的最佳逼近元,则$\forall x \in C, \forall t \in [0,1]$,有$tx+(1-t)x_0\in C$,故
\begin{align*}
\left\| y-\left( tx+(1-t)x_0 \right) \right\| \geqslant \left\| y-x_0 \right\| \Longrightarrow \varphi _x\left( t \right) =\left\| y-\left( tx+(1-t)x_0 \right) \right\| ^2\geqslant \left\| y-x_0 \right\| ^2=\varphi _x\left( 0 \right) .
\end{align*}
又注意到
\begin{align*}
\varphi _x\left( t \right) \geqslant \varphi _x\left( 0 \right) \left( \forall x\in C,t\in \left[ 0,1 \right] \right) \Longrightarrow \varphi _x\left( 1 \right) \geqslant \varphi _x\left( 0 \right)\left( \forall x\in C \right) \Longrightarrow \left\| y-x \right\| \geqslant \left\| y-x_0 \right\| \left( \forall x\in C \right) .
\end{align*}
故$x_0$是$y$在$C$上的最佳逼近元当且仅当
\begin{align}
\varphi_x(t) \geqslant \varphi_x(0) \quad (\forall x \in C, \forall t \in [0,1]). \label{eq:phi_min}
\end{align}
下证\eqref{eq:best_approx_char}式成立$\iff$ \eqref{eq:phi_min}式成立。
因为
\begin{align}\label{eq::JHIOSJDFJIOASIO}
\varphi_x(t) = \|(y - x_0) + t(x_0 - x)\|^2 
= \|y - x_0\|^2 + 2t\operatorname{Re}(y - x_0, x_0 - x) + t^2\|x_0 - x\|^2,
\end{align}
所以
\begin{align*}
\varphi_x'(0) = 2\operatorname{Re}(y - x_0, x_0 - x). 
\end{align*}
因此
\begin{align}
\eqref{eq:best_approx_char}式成立 \iff \varphi_x'(0) \geqslant 0. \label{eq:char_iff_deriv}
\end{align}
又由\eqref{eq::JHIOSJDFJIOASIO}式知
\begin{align*}
\varphi_x(t) - \varphi_x(0) = \varphi_x'(0)t + \|x_0 - x\|^2 t^2,
\end{align*}
所以
\begin{align}
\varphi_x'(0) \geqslant 0 \iff \eqref{eq:phi_min}式成立. \label{eq:deriv_iff_min}
\end{align}
联合\eqref{eq:char_iff_deriv}式与\eqref{eq:deriv_iff_min}式即得结论。
\end{proof}

\begin{corollary}\label{corollary:闭线性流形上的最佳逼近元充要条件}
设$M$是Hilbert空间$\mathscr{X}$的一个闭线性子流形。对$\forall x \in \mathscr{X}$，$y$是$x$在$M$上的最佳逼近元当且仅当它适合
\begin{align}
x - y \perp M - \{y\} \triangleq \{w = z - y \mid \forall z \in M\}. \label{eq:best_approx_linear}
\end{align}
\end{corollary}
\begin{proof}
由\refthe{theorem:泛函分析--最佳逼近元的刻画}，$y$是$x$在$M$上的最佳逼近元当且仅当
\begin{align}
\operatorname{Re}(x - y, y - z) \geqslant 0 \quad (\forall z \in M). \label{eq:linear_char}
\end{align}
因此充分性由上式显然成立,下证必要性.
因为$M$是线性流形，所以$\forall z \in M$可表示为
\begin{align}
z = y + w \quad (w \in M - \{y\}). \label{eq:linear_rep}
\end{align}
注意到$M - \{y\}$是线性子空间，且当$z$跑遍$M$时，$w$跑遍$M - \{y\}$。将\eqref{eq:linear_rep}式代入\eqref{eq:linear_char}式得
\begin{align}
\mathrm{Re}(x-y,-w)\geqslant 0\Longrightarrow \mathrm{Re}(x-y,w)\leqslant 0\quad (\forall w\in M-\{y\}).\label{eq:re_ineq}
\end{align}
在\eqref{eq:re_ineq}式中，用$-w$代替$w$得
\begin{align*}
\mathrm{Re}\left( x-y,-w \right) \leqslant 0\left( \forall -w\in M-\left\{ y \right\} \right) \Longrightarrow \mathrm{Re}(x-y,-w)\geqslant 0\left( \forall w\in M-\left\{ y \right\} \right) \Longrightarrow \mathrm{Re}(x-y,w)\leqslant 0\left( \forall w\in M-\left\{ y \right\} \right) .
\end{align*}
进而
\begin{align}
\operatorname{Re}(x - y, w) = 0 \quad (\forall w \in M - \{y\}). \label{eq:re_eq}
\end{align}
进一步在\eqref{eq:re_eq}式中用$\mathrm{i}w$代替$w$，便推出
\begin{align*}
\mathrm{Re}\left( x-y,\mathrm{i}w \right) =0\left( \forall \mathrm{i}w\in M-\left\{ y \right\} \right) &\Longrightarrow \mathrm{Re}\left( x-y,\mathrm{i}w \right) =0\left( \forall w\in M-\left\{ y \right\} \right) 
\\
&\Longrightarrow -\mathrm{iRe}\left( x-y,w \right) =0\left( \forall w\in M-\left\{ y \right\} \right) 
\\
&\Longrightarrow \mathrm{Im}\left( x-y,w \right) =0\left( \forall w\in M-\left\{ y \right\} \right) .
\end{align*}
故
\begin{align*}
(x - y, w) = 0 \quad (\forall w \in M - \{y\}).
\end{align*}
这就是\eqref{eq:best_approx_linear}式。
\end{proof}

\begin{corollary}\label{corollary:泛函分析--推论289293}
设$M$是Hilbert空间$\mathscr{X}$的一个闭线性子空间,$\forall y\in M$,则$y$是$x$在$M$上的最佳逼近元当且仅当$x - y \perp M$。
\end{corollary}
\begin{proof}
注意到对$\forall y\in M$,都有$M - \{y\} = M$.再由\refcor{corollary:闭线性流形上的最佳逼近元充要条件}立得.
\end{proof}

\begin{corollary}[正交分解]\label{corollary:泛函分析--正交分解定理}
设$M$是Hilbert空间$\mathscr{X}$上的一个闭线性子空间。则对$\forall x \in \mathscr{X}$，存在下列唯一的\textbf{正交分解}：
\begin{align}
x = y + z \quad (y \in M, z \in M^\perp). \label{eq:orthogonal_decomp}
\end{align}
由$x$的正交分解产生的$y$称为$x$在$M$上的\textbf{正交投影}(见\reffig{figure:orthogonal_decomp})。
\end{corollary}
\begin{proof}
由\refcor{corollary:泛函分析--最佳逼近元的存在唯一性}知,可取$y$为$x$在$M$上的最佳逼近元，而$z = x - y$，再由\refcor{corollary:泛函分析--推论289293}知其为满足\eqref{eq:orthogonal_decomp}式的分解。又若还存在另外一种分解：
\begin{align*}
x = y' + z' \quad (y' \in M, z' \in M^\perp),
\end{align*}
则
\begin{align*}
y - y' = z - z' \in (M \cap M^\perp) = \{\theta\},
\end{align*}
即得分解的唯一性.
\end{proof}
\begin{figure}[H]
\centering
\includegraphics[scale=0.3]{泛函分析--图1.6.5.png}
\caption{}
\label{figure:orthogonal_decomp}
\end{figure}


\subsection{应用:最小二乘法}

\begin{example}\label{example:泛函分析--最小二乘法}
\begin{enumerate}[(1)]
\item \textbf{实际观测问题}. 已知量 $y$ 与量 $x_1, x_2, \cdots, x_n$ 之间呈线性关系:
\begin{align*}
y = \lambda_1 x_1 + \lambda_2 x_2 + \cdots + \lambda_n x_n.
\end{align*}
为确定未知系数 $\lambda_1, \lambda_2, \cdots, \lambda_n$，观测 $m$ 组数据（$m > n$），即
\begin{table}[H]
\centering
\caption{}
\begin{tabular}{ccccc}
$y^{(1)},$ & $x_1^{(1)},$ & $x_2^{(1)},$ & $\cdots,$ & $x_n^{(1)}$ \\
$y^{(2)},$ & $x_1^{(2)},$ & $x_2^{(2)},$ & $\cdots,$ & $x_n^{(2)}$ \\
$y^{(3)},$ & $x_1^{(3)},$ & $x_2^{(3)},$ & $\cdots,$ & $x_n^{(3)}$ \\
$\vdots$   & $\vdots$     & $\vdots$     & $\ddots$  & $\vdots$   \\
$y^{(m)},$ & $x_1^{(m)},$ & $x_2^{(m)},$ & $\cdots,$ & $x_n^{(m)}$ \\
\end{tabular}
\end{table}
求系数 $\lambda_1, \lambda_2, \cdots, \lambda_n$，使得
\begin{align*}
\min_{\alpha_1, \alpha_2, \cdots, \alpha_n} \sum\limits_{j=1}^{m} \left| y^{(j)} - \sum\limits_{i=1}^{n} \alpha_i x_i^{(j)} \right|^2 = \sum\limits_{j=1}^{m} \left| y^{(j)} - \sum\limits_{i=1}^{n} \lambda_i x_i^{(j)} \right|^2.
\end{align*}
该问题可视为在空间 $\mathbb{R}^m$ 中，求由数据表中后 $n$ 列向量张成的子空间上的元素，对第一列向量的最佳逼近.

\item \textbf{平方平均逼近}. 在函数逼近论中，对 $f \in L^2[a,b]$，用给定的 $n$ 个 $L^2[a,b]$ 函数 $\varphi_1, \varphi_2, \cdots, \varphi_n$ 的线性组合按平方平均意义求最佳逼近，即求系数 $\lambda_1, \lambda_2, \cdots, \lambda_n$，使得
\begin{align*}
\min_{\alpha_1, \alpha_2, \cdots, \alpha_n} \int_{a}^{b} \left| f(x) - \sum\limits_{i=1}^{n} \alpha_i \varphi_i(x) \right|^2 \mathrm{d}x = \int_{a}^{b} \left| f(x) - \sum\limits_{i=1}^{n} \lambda_i \varphi_i(x) \right|^2 \mathrm{d}x.
\end{align*}

\item \textbf{最佳估计问题}. 设 $(\Omega, \mathscr{B}, P)$ 是概率空间，$X, X_1, X_2, \cdots, X_n \in L^2(\Omega, \mathscr{B}, P)$，用 $X_1, X_2, \cdots, X_n$ 的线性组合估计 $X$，求系数 $\lambda_1, \lambda_2, \cdots, \lambda_n$，使得
\begin{align*}
\min_{(\alpha_1, \alpha_2, \cdots, \alpha_n) \in \mathbb{R}^n} \int_{\Omega} \left| X(\omega) - \sum\limits_{i=1}^{n} \alpha_i X_i(\omega) \right|^2 \mathrm{d}P(\omega) = \int_{\Omega} \left| X(\omega) - \sum\limits_{i=1}^{n} \lambda_i X_i(\omega) \right|^2 \mathrm{d}P(\omega).
\end{align*}
\end{enumerate}
\end{example}

\begin{theorem}
设 $\mathscr{X}$ 为 Hilbert 空间，$x_1, x_2, \cdots, x_n \in \mathscr{X}$ 线性无关，记 $\mathscr{M} = \operatorname{span}\{x_1, x_2, \cdots, x_n\}$，则存在唯一的 $x_0 = \sum\limits_{i=1}^{n} \lambda_i x_i \in \mathscr{M}$，使得
\begin{align*}
\left\| x - x_0 \right\| = \inf_{y \in \mathscr{M}} \left\| x - y \right\|,
\end{align*}
且
\begin{align*}
\lambda_i = \frac{
\begin{vmatrix}
(x_1, x_1) & \cdots & (x, x_1) & \cdots & (x_n, x_1) \\
(x_1, x_2) & \cdots & (x, x_2) & \cdots & (x_n, x_2) \\
\vdots & & \vdots & & \vdots \\
(x_1, x_n) & \cdots & (x, x_n) & \cdots & (x_n, x_n)
\end{vmatrix}
}{
\begin{vmatrix}
(x_1, x_1) & \cdots & (x_i, x_1) & \cdots & (x_n, x_1) \\
(x_1, x_2) & \cdots & (x_i, x_2) & \cdots & (x_n, x_2) \\
\vdots & & \vdots & & \vdots \\
(x_1, x_n) & \cdots & (x_i, x_n) & \cdots & (x_n, x_n)
\end{vmatrix}
}, \quad i = 1, 2, \cdots, n.
\end{align*}
\end{theorem}
\begin{remark}
\refexa{example:泛函分析--最小二乘法}中的三个问题本质上等价于这个定理.
\end{remark}
\begin{proof}
设 $x_0 = \sum\limits_{i=1}^{n} \lambda_i x_i$ 为所求元素，则由\refcor{corollary:泛函分析--推论289293}知其为 $x$ 在 $\mathscr{M}$ 上的正交投影当且仅当 $x - x_0 \perp \mathscr{M}$，即对任意 $j = 1, 2, \cdots, n$，有 $(x - x_0, x_j) = 0$.

将 $x_0 = \sum\limits_{i=1}^{n} \lambda_i x_i$ 代入内积条件，得
\begin{align*}
\left( x - \sum\limits_{i=1}^{n} \lambda_i x_i, x_j \right) = 0, \quad j = 1, 2, \cdots, n,
\end{align*}
展开即得
\begin{align}
\sum\limits_{i=1}^{n} \lambda_i (x_i, x_j) = (x, x_j), \quad j = 1, 2, \cdots, n. \label{eq:1_6_34}
\end{align}

由于 $x_1, x_2, \cdots, x_n$ 线性无关，Gram 矩阵 $((x_i, x_j))_{n \times n}$ 可逆，故方程组 \eqref{eq:1_6_34} 存在唯一解 $(\lambda_1, \lambda_2, \cdots, \lambda_n) \in \mathbb{R}^n$，即 $x_0$ 存在唯一.由\hyperref[Basis of Algebra-theorem:Crammer法则]{Crammer法则}可得
\begin{align*}
\lambda_i = \frac{
\begin{vmatrix}
(x_1, x_1) & \cdots & (x, x_1) & \cdots & (x_n, x_1) \\
(x_1, x_2) & \cdots & (x, x_2) & \cdots & (x_n, x_2) \\
\vdots & & \vdots & & \vdots \\
(x_1, x_n) & \cdots & (x, x_n) & \cdots & (x_n, x_n)
\end{vmatrix}
}{
\begin{vmatrix}
(x_1, x_1) & \cdots & (x_i, x_1) & \cdots & (x_n, x_1) \\
(x_1, x_2) & \cdots & (x_i, x_2) & \cdots & (x_n, x_2) \\
\vdots & & \vdots & & \vdots \\
(x_1, x_n) & \cdots & (x_i, x_n) & \cdots & (x_n, x_n)
\end{vmatrix}
}, \quad i = 1, 2, \cdots, n.
\end{align*}
\end{proof}


\begin{proposition}[极化恒等式]
设 $a$ 是复线性空间 $\mathscr{X}$ 上的共轭双线性函数, $q$ 是由 $a$ 诱导的二次型, 求证: 对 $\forall x,y\in\mathscr{X}$ 有
\begin{align*}
a(x,y)=\frac{1}{4}\left[q(x+y)-q(x-y)+\mathrm{i}q(x+\mathrm{i}y)-\mathrm{i}q(x-\mathrm{i}y)\right].
\end{align*}
\end{proposition}

\begin{example}
求证: 在 $C[a,b]$ 中不可能引进一种内积 $(\cdot,\cdot)$, 使其满足
\begin{align*}
(f,f)^\frac{1}{2}=\max_{a\leqslant x\leqslant b}|f(x)|\quad(\forall f\in C[a,b]).
\end{align*}
\end{example}

\begin{example}
在 $L^2[0,T]$ 中, 求证: 函数
\begin{align*}
x\mapsto\left|\int_0^T\mathrm{e}^{-(T-\tau)}x(\tau)\mathrm{d}\tau\right|\quad(\forall x\in L^2[0,T])
\end{align*}
在单位球面上达到最大值, 并求出此最大值和达到最大值的元素 $x$.
\end{example}
\begin{note}
提示 利用 Cauchy-Schwarz 不等式及其取等号的条件.
\end{note}

\begin{example}
在 $L^2[-1,1]$ 中, 问偶函数集的正交补是什么? 证明你的结论.
\end{example}

\begin{example}
在 $L^2[a,b]$ 中, 考察函数集 $S=\{\mathrm{e}^{2\pi \mathrm{i}nx}\}$.
\begin{enumerate}[(1)]
\item 若 $|b-a|\leqslant 1$, 求证: $S^\perp=\{\theta\}$.
\item 若 $|b-a|>1$, 求证: $S^\perp\neq\{\theta\}$.
\end{enumerate}
\end{example}

\begin{example}
设 $\mathscr{X}$ 表示闭单位圆上的解析函数全体, 内积定义为
\begin{align*}
(f,g)=\frac{1}{i}\oint_{|z|=1}\frac{f(z)\overline{g(z)}}{z}\mathrm{d}z\quad(\forall f,g\in\mathscr{X}).
\end{align*}
求证: $\{z^n/\sqrt{2\pi}\}$ 是一组正交规范集.
\end{example}

\begin{example}
设 $\{e_n\}_{n=1}^\infty,\{f_n\}_{n=1}^\infty$ 是 Hilbert 空间 $\mathscr{X}$ 中的两个正交规范集, 满足条件
\begin{align*}
\sum\limits_{n=1}^\infty\|e_n-f_n\|^2<1.
\end{align*}
求证: $\{e_n\}$ 和 $\{f_n\}$ 两者中一个完备蕴含另一个完备.
\end{example}

\begin{example}
设 $\mathscr{X}$ 是 Hilbert 空间, $\mathscr{X}_0$ 是 $\mathscr{X}$ 的闭线性子空间, $\{e_n\},\{f_n\}$ 分别是 $\mathscr{X}_0$ 和 $\mathscr{X}_0^\perp$ 的正交规范基. 求证: $\{e_n\}\cup\{f_n\}$ 是 $\mathscr{X}$ 的正交规范基.
\end{example}

\begin{example}
设 $H^2(D)$ 是按例定义的内积空间.
\begin{enumerate}[(1)]
\item 如果 $u(z)$ 的 Taylor 展开式是 $u(z)=\sum\limits_{k=0}^\infty b_k z^k$, 求证:
\begin{align*}
\sum\limits_{k=0}^\infty\frac{|b_k|^2}{1+k}<\infty.
\end{align*}
\item 设 $u(z),v(z)\in H^2(D)$, 并且
\begin{align*}
u(z)=\sum\limits_{k=0}^\infty a_k z^k,\quad v(z)=\sum\limits_{k=0}^\infty b_k z^k,
\end{align*}
求证:
\begin{align*}
(u,v)=\pi\sum\limits_{k=0}^\infty\frac{a_k\overline{b_k}}{k+1}.
\end{align*}
\item 设 $u(z)\in H^2(D)$, 求证:
\begin{align*}
|u(z)|\leqslant\frac{\|u\|}{\sqrt{\pi}(1-|z|)}\quad(\forall|z|<1).
\end{align*}
\item 验证 $H^2(D)$ 是 Hilbert 空间.
\end{enumerate}
\end{example}

\begin{example}
设 $\mathscr{X}$ 是内积空间, $\{e_n\}$ 是 $\mathscr{X}$ 中的正交规范集, 求证:
\begin{align*}
\left|\sum\limits_{n=1}^\infty(x,e_n)\overline{(y,e_n)}\right|\leqslant\|x\|\cdot\|y\|\quad(\forall x,y\in\mathscr{X}).
\end{align*}
\end{example}

\begin{example}
设 $\mathscr{X}$ 是一个内积空间, $\forall x_0\in\mathscr{X},\forall r>0$, 令
\begin{align*}
C=\{x\in\mathscr{X}\mid\|x-x_0\|\leqslant r\}.
\end{align*}
\begin{enumerate}[(1)]
\item 求证: $C$ 是 $\mathscr{X}$ 中的闭凸集.
\item $\forall x\in\mathscr{X}$, 令
\begin{align*}
y=
\begin{cases}
x_0+r\frac{x-x_0}{\|x-x_0\|}, & \text{当}\ x\notin C, \\
x, & \text{当}\ x\in C.
\end{cases}
\end{align*}
求证: $y$ 是 $x$ 在 $C$ 中的最佳逼近元.
\end{enumerate}
\end{example}

\begin{example}
求 $(a_0,a_1,a_2)\in\mathbb{R}^3$, 使得
\begin{align*}
\int_0^1|\mathrm{e}^t-a_0-a_1t-a_2t^2|^2\mathrm{d}t
\end{align*}
取最小值.
\end{example}

\begin{example}
设 $f(x)\in C^2[a,b]$, 满足边界条件
\begin{align*}
f(a)=f(b)=0,\quad f'(a)=1,\quad f'(b)=0.
\end{align*}
求证:
\begin{align*}
\int_a^b|f''(x)|^2\mathrm{d}x\geqslant\frac{4}{b-a}.
\end{align*}
\end{example}

\begin{proposition}[变分不等式]
设 $\mathscr{X}$ 是一个 Hilbert 空间, $a(x,y)$ 是 $\mathscr{X}$ 上的共轭对称的双线性函数, $\exists M>0,\delta>0$, 使得
\begin{align*}
\delta\|x\|^2\leqslant a(x,x)\leqslant M\|x\|^2\quad(\forall x\in\mathscr{X}).
\end{align*}
又设 $u_0\in\mathscr{X}$, $C$ 是 $\mathscr{X}$ 上的一个闭凸子集. 求证: 函数
\begin{align*}
x\mapsto a(x,x)-\text{Re}(u_0,x)
\end{align*}
在 $C$ 上达到最小值, 并且达到最小值的点 $x_0$ 唯一, 且满足
\begin{align*}
\text{Re}\left[2a(x_0,x-x_0)-(u_0,x-x_0)\right]\geqslant 0\quad(\forall x\in C).
\end{align*}
\end{proposition}









\end{document}